% main.tex -- template for journal article, using American Physical Society / American Institute of Physics format
% -- see https://www.overleaf.com/latex/templates/template-for-submission-to-aip-publishing-journals/xhsskcchtbxf
% ~ Mark J. Duvall ~ duvall3.git@gmail.com ~ 05/2022 ~ %

%Copyright (C) 2022 Mark J. Duvall / T. Rocks Science
%
%    This program is free software: you can redistribute it and/or modify
%    it under the terms of the GNU General Public License as published by
%    the Free Software Foundation, either version 3 of the License, or
%    (at your option) any later version.
%
%    This program is distributed in the hope that it will be useful,
%    but WITHOUT ANY WARRANTY; without even the implied warranty of
%    MERCHANTABILITY or FITNESS FOR A PARTICULAR PURPOSE.  See the
%    GNU General Public License for more details.
%
%    You should have received a copy of the GNU General Public License
%    along with this program.  If not, see <http://www.gnu.org/licenses/>.


%% preamble
\documentclass[twocolumn]{revtex4-2}	% APS/AIP article
\usepackage[utf8]{inputenc}		% input encoding
\usepackage{graphicx}			% use graphics
\usepackage{lipsum}			% for placeholder text
% for hyperlinks:
\usepackage{hyperref}
\hypersetup{colorlinks=true, linkcolor=blue, urlcolor=blue, citecolor=blue}
% for tables:
\usepackage{tabularx}
\usepackage{multirow}
\usepackage{array}
% for figures:
\usepackage{caption}
\usepackage{subcaption}
\usepackage[percent]{overpic}
% useful packages for a physics paper:
\usepackage{epsfig}			% rendering certain image files
\usepackage{xcolor}			% common colors
\usepackage{amsmath}			% common math symbols
\usepackage{amssymb}			% additional math symbols
\usepackage{bm}				% bold math symbols
\usepackage{units}			% common units


% other preamble items like definitions can go here; for example,
\definecolor{darkgreen}{RGB}{0,175,0}


%% begin document
\begin{document}


% acronyms can go here; for example,
\newcommand{\ke}{kinetic energy (KE)\renewcommand{\ke}{KE}}
\newcommand{\pe}{potential energy (PE)\renewcommand{\pe}{PE}}

% other useful document-wide commands can go here as well; for example,
\newcommand{\dash}{\textbf{---}\ } % command to produce an em-dash that is actually heavy enough to see
\newcommand{\bs}{\boldsymbol} % math-mode shortcut for bold symbols -- e.g., $ \bs{r} \equiv \langle x, y, z \rangle $
\newcommand{\nucl}[2]{{}^{#1}\mbox{#2}} % math-mode command for writing isotopes -- e.g., write lithium-6 using $\nucl{6}{Li}$

% title
\title{TITLE}

% author(s)
\author{A.~B.~LastnameOne} % F(irst initial) M(iddle initial) Lastname
\email{author1@domain.suffix}
\affiliation{Department, Institution, City, Country}
\author{C.~D.~LastnameTwo}
\author{E.~F.~LastnameThree}
\author{G.~H.~LastnameFour}

% date
\date{\today}

% abstract
\begin{abstract}
\emph{\textbf{Example Abstract text.}}
\lipsum[91]
\end{abstract}

% make title and TOC
\maketitle %\maketitle must follow \title, \authors, \abstract and \pacs, as applicable
\tableofcontents


%% MAIN -- sections and bibliography
\section{\label{sec:intro}Introduction}

\emph{\textbf{Example Introduction text.}}
\lipsum[1]

\begin{equation}
  \bs{r} \equiv \langle x, y, z \rangle
  \label{eq:r_defn}
\end{equation}

Example reference to Eq.~\ref{eq:r_defn}.\\

\lipsum[2]

\section{\label{sec:body}Body}

% example figure
\begin{figure}[ht]
  \centering
  \includegraphics[width=0.5\linewidth]{figures/seal.jpg}
  \caption{Caption for Example Figure}
  \label{fig:exmp}
\end{figure}

\emph{\textbf{Example Body text.}}
Example citation~\cite{Apollonio:1999jg}, example reference to Fig.~\ref{fig:exmp}.
Example expansion of acronym for \ke.
Example later usage of \ke.
Example reference to~\S\ref{sec:concl}.
Eample usage of user-defined \texttt{\textbackslash nucl} command: $\nucl{6}{Li}$.
\\

\lipsum[4-8]

\section{\label{sec:concl}Conclusion}

\emph{\textbf{Example Conclusion text.}}
\lipsum[9-10]

\bibliography{Refs}\bibliographystyle{aipnum4-2} % aip-style bibliography


%% end document
\end{document}

%% all pau!   )
